\documentclass[12pt]{article}

% --- PACKAGES ---
\usepackage[utf8]{inputenc}
\usepackage[a4paper, margin=1in]{geometry}
\usepackage{graphicx}
\usepackage{amsmath}
\usepackage{booktabs}
\usepackage{hyperref}
\usepackage{float}
\usepackage{enumitem}
\usepackage{array}
\usepackage{listings}
\usepackage{xcolor}
\usepackage{mwe} % For placeholder images, you can remove this

% --- HYPERLINK AND LISTING STYLES ---
\hypersetup{
    colorlinks=true,
    linkcolor=blue,
    filecolor=magenta,      
    urlcolor=cyan,
}

\definecolor{codegreen}{rgb}{0,0.6,0}
\definecolor{codegray}{rgb}{0.5,0.5,0.5}
\definecolor{codepurple}{rgb}{0.58,0,0.82}
\definecolor{backcolour}{rgb}{0.95,0.95,0.92}

\lstdefinestyle{mystyle}{
    backgroundcolor=\color{backcolour},   
    commentstyle=\color{codegreen},
    keywordstyle=\color{magenta},
    numberstyle=\tiny\color{codegray},
    stringstyle=\color{codepurple},
    basicstyle=\footnotesize\ttfamily,
    breakatwhitespace=false,         
    breaklines=true,                 
    captionpos=b,                    
    keepspaces=true,                 
    numbers=left,                    
    numbersep=5pt,                  
    showspaces=false,                
    showstringspaces=false,
    showtabs=false,                  
    tabsize=2
}
\lstset{style=mystyle}

%====================================================================
% DOCUMENT START
%====================================================================
\begin{document}

% --- TITLE PAGE ---
\begin{center}
    \large \textbf{Sri Sivasubramaniya Nadar College of Engineering, Chennai} \\
    (An autonomous Institution affiliated to Anna University) \\
    \vspace{1cm}
\end{center}

\begin{table}[H]
\renewcommand{\arraystretch}{1.5}
\resizebox{\textwidth}{!}{%
\begin{tabular}{|l|cll|}
\hline
Degree \& Branch      & \multicolumn{1}{c|}{B.E. Computer Science \& Engineering} & \multicolumn{1}{l|}{Semester}         & V                     \\ \hline
Subject Code \& Name & \multicolumn{3}{c|}{ICS1512 - Machine Learning Algorithms Laboratory} \\ \hline
Academic year         & \multicolumn{1}{c|}{2025-2026 (Odd)}                      & \multicolumn{1}{c|}{Batch: 2023-2028} & \multicolumn{1}{c|}{\textbf{ Date: 12/08/2025}} \\ \hline
\end{tabular}%
}
\end{table}

\vspace{1cm}

\begin{center}
 \section*{Experiment 3: Spam Mail Prediction using Multiple Classification Models}
\end{center}

\vspace{1cm}

\section{Aim}
To build, evaluate, and compare multiple classification models for predicting whether an email is spam or not, justifying the model choices and using cross-validation to identify the most suitable model.

\section{Libraries Used}
Based on the import statements in the notebook, the primary libraries used were:
\begin{itemize}
    \item \textbf{pandas} – For data manipulation and reading CSV files.
    \item \textbf{numpy} – For numerical computations.
    \item \textbf{matplotlib} – For creating visualizations like plots.
    \item \textbf{seaborn} – For statistical data visualization.
    \item \textbf{scikit-learn} – For implementing and evaluating machine learning models, including metrics, model selection, and preprocessing.
    \item \textbf{xgboost} – For implementing the XGBoost classifier.
\end{itemize}

\section{Objectives}
\begin{itemize}
    \item To apply and compare a wide range of classification algorithms for a predictive modeling task.
    \item To evaluate model performance using Accuracy, Precision, Recall, and F1-Score.
    \item To implement cross-validation to ensure model performance is robust and generalizable.
    \item To identify the best-performing model for the given dataset and problem.
\end{itemize}

%====================================================================
% PREPROCESSING SECTION
%====================================================================

\section{Data Preprocessing And EDA }
\begin{verbatim}
#Data Loading
import pandas as pd
import numpy as np
import matplotlib.pyplot as plt
import seaborn as sns

# Load the data
df = pd.read_csv('/content/drive/MyDrive/ML_LAB/mail.csv')
# Display the first 5 rows
print("First 5 rows of the dataset:")
print(df.head())
print("Initial shape:", df.shape)
\end{verbatim}

\begin{figure}[H]
    \centering
    \includegraphics[width=0.6\textwidth]{dataset.png}
    \caption{Output , showing the first five records of the dataset.}
    \label{fig:dataset_head}
\end{figure}


\subsection*{Handling Missing Values and Outliers}
\begin{verbatim}
# HANDLE MISSING VALUES
df = df.dropna(thresh=df.shape[1]//2)  # Drop rows with >50% missing
df.fillna(df.median(numeric_only=True), inplace=True)

# OUTLIER HANDLING (Z-Score)
def remove_outliers(df, threshold=3):
    numeric_cols = df.select_dtypes(include=[np.number]).columns
    z_scores = np.abs((df[numeric_cols] - df[numeric_cols].mean()) / df[numeric_cols].std())
    return df[(z_scores < threshold).all(axis=1)]
df = remove_outliers(df)
print("After outlier removal:", df.shape)

After outlier removal: (2185, 58)

# FEATURE / TARGET SPLIT
X = df.drop(columns=[TARGET_COLUMN])
y = df[TARGET_COLUMN]

# ENCODE + STANDARDIZE
numeric_cols = X.select_dtypes(include=[np.number]).columns
categorical_cols = X.select_dtypes(exclude=[np.number]).columns
X_encoded = pd.get_dummies(X, columns=categorical_cols)
scaler = StandardScaler()
X_scaled = scaler.fit_transform(X_encoded)

\end{verbatim}

%====================================================================
% MODELING SECTION
%====================================================================
\section{Model Training and Evaluation}

\subsection*{Train-Validation-Test Split}
The data was partitioned into training (70\%), validation (15\%), and test (15\%) sets to ensure robust model evaluation.

\begin{verbatim}
# TRAIN / VAL / TEST SPLIT
X_train, X_temp, y_train, y_temp = train_test_split(X_encoded, y, test_size=0.3, random_state=42)
X_val, X_test, y_val, y_test = train_test_split(X_temp, y_temp, test_size=0.5, random_state=42)

print("Train set:", X_train.shape)
print("Validation set:", X_val.shape)
print("Test set:", X_test.shape)
\end{verbatim}

\begin{verbatim}
Train set: (1529, 57)
Validation set: (328, 57)
Test set: (328, 57)
\end{verbatim}
\subsection*{Model Training}
\begin{verbatim}
def evaluate_model(name, model, X_train, X_test, param_grid=None, param_dist=None):

    print(f"\n{name} - Hyperparameter Tuning Started")

    #GRIDSEARCHCV
    if param_grid:
        grid_search = GridSearchCV(
            estimator=model,
            param_grid=param_grid,
            cv=5,
            scoring='accuracy',
            n_jobs=-1
        )
        grid_search.fit(X_train, y_train)
        print(f"Best Params (GridSearchCV): {grid_search.best_params_}")
        print(f"Best CV Score (GridSearchCV): {grid_search.best_score_:.4f}")
    else:
        grid_search = None

    # RANDOMIZEDSEARCHCV
    if param_dist:
        random_search = RandomizedSearchCV(
            estimator=model,
            param_distributions=param_dist,
            n_iter=10,
            cv=5,
            scoring='accuracy',
            random_state=42,
            n_jobs=-1
        )
        random_search.fit(X_train, y_train)
        print(f"Best Params (RandomizedSearchCV): {random_search.best_params_}")
        print(f"Best CV Score (RandomizedSearchCV): {random_search.best_score_:.4f}")
    else:
        random_search = None

    # PICK BEST MODEL
    if grid_search and random_search:
        best_model = (
            grid_search if grid_search.best_score_ >= random_search.best_score_
            else random_search
        ).best_estimator_
    elif grid_search:
        best_model = grid_search.best_estimator_
    elif random_search:
        best_model = random_search.best_estimator_
    else:
        best_model = model

    #  EVALUATION
    start_time = time.time()
    best_model.fit(X_train, y_train)
    end_time = time.time()

    y_pred = best_model.predict(X_test)
    print(f"\n{name} Performance:")
    print(f"Accuracy :  {accuracy_score(y_test, y_pred):.4f}")
    print(f"Precision:  {precision_score(y_test, y_pred, average='macro'):.4f}")
    print(f"Recall   :  {recall_score(y_test, y_pred, average='macro'):.4f}")
    print(f"F1 Score :  {f1_score(y_test, y_pred, average='macro'):.4f}")
    print(f"Training Time: {(end_time - start_time):.4f} seconds")

    fig, axes = plt.subplots(1, 2, figsize=(12, 5))

    cm = confusion_matrix(y_test, y_pred)
    disp = ConfusionMatrixDisplay(confusion_matrix=cm)
    disp.plot(ax=axes[0], cmap='Blues')
    axes[0].set_title('Confusion Matrix')

    if hasattr(best_model, "predict_proba"):
        y_prob = best_model.predict_proba(X_test)[:, 1]
        fpr, tpr, _ = roc_curve(y_test, y_prob)
        roc_auc = auc(fpr, tpr)

        axes[1].plot(fpr, tpr, label=f'AUC = {roc_auc:.2f}', color='green')
        axes[1].plot([0, 1], [0, 1], linestyle='--', color='blue')
        axes[1].set_xlabel('False Positive Rate')
        axes[1].set_ylabel('True Positive Rate')
        axes[1].set_title('ROC Curve')
        axes[1].legend()
        axes[1].grid(True)
    else:
        axes[1].axis('off')

    plt.tight_layout()
    plt.show()
\end{verbatim}
\subsection*{Model Evaluation Results}
The following models were trained and evaluated.
\subsubsection*{Logistic Regression}
\begin{verbatim}
evaluate_model("Logistic Regression", LogisticRegression(max_iter=1000), X_train, X_test)
Logistic Regression Performance:
Accuracy :  0.9207
Precision:  0.9287
Recall   :  0.9098
F1 Score :  0.9167
Training Time: 3.7717 seconds
\end{verbatim}
\begin{figure}[H]
    \centering
    \includegraphics[width=\textwidth]{logistic_Regression.png}
    \caption{Confusion Matrix and ROC Curve for Logistic Regression}
\end{figure}

\subsubsection*{Naive Bayes - Gaussian}
\begin{verbatim}
evaluate_model("GaussianNB", GaussianNB(), X_train, X_test)
GaussianNB Performance:
Accuracy :  0.8140
Precision:  0.8304
Recall   :  0.8347
F1 Score :  0.8139
Training Time: 0.0058 seconds
\end{verbatim}
\begin{figure}[H]
    \centering
    \includegraphics[width=\textwidth]{GaussianNB.png}
    \caption{Confusion Matrix and ROC Curve for GaussianNB}
\end{figure}

\subsubsection*{Naive Bayes - Multinomial}
\begin{verbatim}
evaluate_model("MultinomialNB", MultinomialNB(), X_train, X_test)
MultinomialNB Performance:
Accuracy :  0.7744
Precision:  0.7677
Recall   :  0.7665
F1 Score :  0.7671
Training Time: 0.0053 seconds
\end{verbatim}
\begin{figure}[H]
    \centering
    \includegraphics[width=\textwidth]{MultinomialNB.png}
    \caption{Confusion Matrix and ROC Curve for MultinomialNB}
\end{figure}

\subsubsection*{Naive Bayes - Bernoulli}
\begin{verbatim}
evaluate_model("BernoulliNB", BernoulliNB(), X_train, X_test)
BernoulliNB Performance:
Accuracy :  0.8811
Precision:  0.8837
Recall   :  0.8706
F1 Score :  0.8756
Training Time: 0.0068 seconds
\end{verbatim}
\begin{figure}[H]
    \centering
    \includegraphics[width=0.9\textwidth]{BernoulliNB.png}
    \caption{Confusion Matrix and ROC Curve for BernoulliNB}
\end{figure}

\subsubsection*{K-Nearest Neighbors - Varying k values [1, 3, 5, 7, 9]}
\begin{verbatim}
for k in [1, 3, 5, 7, 9]:
    evaluate_model(f"KNN (k={k})", KNeighborsClassifier(n_neighbors=k), X_train, X_test)

KNN (k=1) Performance:
Accuracy :  0.7683
Precision:  0.7616
Recall   :  0.7592
F1 Score :  0.7603
Training Time: 0.0041 seconds
\end{verbatim}
\begin{figure}[H]
    \centering
    \includegraphics[width=0.9\textwidth]{KNN-1.png}
    \caption{Confusion Matrix and ROC Curve for KNN-1}
\end{figure}


\begin{verbatim}
KNN (k=3) Performance:
Accuracy :  0.7591
Precision:  0.7524
Recall   :  0.7482
F1 Score :  0.7499
Training Time: 0.0038 seconds
\end{verbatim}
\begin{figure}[H]
    \centering
    \includegraphics[width=0.9\textwidth]{KNN-3.png}
    \caption{Confusion Matrix and ROC Curve for KNN-3}
\end{figure}
\begin{verbatim}
KNN (k=5) Performance:
Accuracy :  0.7622
Precision:  0.7572
Recall   :  0.7475
F1 Score :  0.7508
Training Time: 0.0036 seconds
\end{verbatim}
\begin{figure}[H]
    \centering
    \includegraphics[width=0.9\textwidth]{KNN-5.png}
    \caption{Confusion Matrix and ROC Curve for KNN-5}
\end{figure}

\begin{verbatim}
KNN (k=7) Performance:
Accuracy :  0.7561
Precision:  0.7528
Recall   :  0.7381
F1 Score :  0.7423
Training Time: 0.0038 seconds
\end{verbatim}
\begin{figure}[H]
    \centering
    \includegraphics[width=0.9\textwidth]{KNN-7.png}
    \caption{Confusion Matrix and ROC Curve for KNN-7}
\end{figure}

\begin{verbatim}
KNN (k=9) Performance:
Accuracy :  0.7348
Precision:  0.7322
Recall   :  0.7123
F1 Score :  0.7166
Training Time: 0.0035 seconds
\end{verbatim}
\begin{figure}[H]
    \centering
    \includegraphics[width=0.9\textwidth]{KNN-9.png}
    \caption{Confusion Matrix and ROC Curve for KNN-9}
\end{figure}

\subsection*{K-Nearest Neighbors - KDTree}
\begin{verbatim}
    evaluate_model("KNN (KDTree)", KNeighborsClassifier(algorithm='kd_tree'), X_train, X_test)
    
KNN (KDTree) Performance:
Accuracy :  0.7622
Precision:  0.7572
Recall   :  0.7475
F1 Score :  0.7508
Training Time: 0.0149 seconds
\end{verbatim}
\begin{figure}[H]
    \centering
    \includegraphics[width=0.9\textwidth]{KNN_KD.png}
    \caption{Confusion Matrix and ROC Curve for KNN-KDTree }
\end{figure}

\subsection*{K-Nearest Neighbors - BallTree}
\begin{verbatim}
    evaluate_model("KNN (BallTree)", KNeighborsClassifier(algorithm='ball_tree'), X_train, X_test)
    
KNN (BallTree) Performance:
Accuracy :  0.7622
Precision:  0.7572
Recall   :  0.7475
F1 Score :  0.7508
Training Time: 0.0092 seconds
\end{verbatim}
\begin{figure}[H]
    \centering
    \includegraphics[width=0.9\textwidth]{KNN_Ball.png}
    \caption{Confusion Matrix and ROC Curve for KNN-BallTree}
\end{figure}

\subsection*{Support Vector Classifier}
\subsubsection*{SVC- Linear kernel}
\begin{verbatim}
svc_linear = SVC(kernel='linear', C=1.0, probability=True)
evaluate_model("SVC (Linear)", svc_linear, X_train, X_test)

SVC (Linear) Performance:
Accuracy :  0.9268
Precision:  0.9353
Recall   :  0.9161
F1 Score :  0.9231
Training Time: 354.6113 seconds
\end{verbatim}
\begin{figure}[H]
    \centering
    \includegraphics[width=0.9\textwidth]{svc_linear.png}
    \caption{Confusion Matrix and ROC Curve for SVC Linear Kernel}
\end{figure}

\subsubsection*{SVC- Polynomial kernel}
\begin{verbatim}
svc_poly = SVC(kernel='poly', C=1.0, degree=3, gamma='scale',probability=True)
evaluate_model("SVC (Poly)", svc_poly, X_train, X_test)

SVC (Poly) Performance:
Accuracy :  0.6524
Precision:  0.8137
Recall   :  0.5809
F1 Score :  0.5248
Training Time: 1.2004 seconds
\end{verbatim}
\begin{figure}[H]
    \centering
    \includegraphics[width=0.9\textwidth]{svc_polynomial.png}
    \caption{Confusion Matrix and ROC Curve for SVC Polynomial Kernel}
\end{figure}

\subsubsection*{SVC- RBF kernel}
\begin{verbatim}
svc_model = SVC(kernel='rbf', C=1.0, gamma='scale', probability=True)
evaluate_model("SVC (RBF Kernel)", svc_model, X_train, X_test)

SVC (RBF Kernel) Performance:
Accuracy :  0.7165
Precision:  0.7742
Recall   :  0.6667
F1 Score :  0.6607
Training Time: 0.6475 seconds
\end{verbatim}
\begin{figure}[H]
    \centering
    \includegraphics[width=0.9\textwidth]{svc_rbf.png}
    \caption{Confusion Matrix and ROC Curve for SVC RBF Kernel}
\end{figure}

\subsubsection*{SVC- Sigmoid kernel}
\begin{verbatim}
svc_sigmoid = SVC(kernel='sigmoid', C=1.0, gamma='scale', probability=True)
evaluate_model("SVC (Sigmoid)", svc_sigmoid, X_train, X_test)

SVC (Sigmoid) Performance:
Accuracy :  0.5366
Precision:  0.5178
Recall   :  0.5173
F1 Score :  0.5170
Training Time: 0.6506 seconds
\end{verbatim}
\begin{figure}[H]
    \centering
    \includegraphics[width=0.9\textwidth]{svc_sigmoid.png}
    \caption{Confusion Matrix and ROC Curve for SVC Sigmoid Kernel}
\end{figure}

\section{Model Justification and Dataset Suitability}
\begin{itemize}
    \item \textbf{Logistic Regression:} Chosen as a simple and interpretable baseline model for this binary classification task.
    \item \textbf{Naive Bayes (Gaussian, Multinomial, Bernoulli):} These probabilistic models are efficient and often perform well on text classification tasks like spam detection, making them a suitable choice.
    \item \textbf{K-Nearest Neighbors:} A non-parametric model chosen to capture potentially complex, non-linear decision boundaries based on feature similarity.
    \item \textbf{Support Vector Machine (SVM):} A powerful model that is effective in high-dimensional spaces. Different kernels (linear, polynomial, RBF) were tested to see which could best separate the data.
\end{itemize}
\section{Model Comparison}
\subsubsection*{Performance Comparison of Na¨ıve Bayes Variants}
\begin{table}[h!]
\centering
\caption{Performance Comparison of Na¨ıve Bayes Variants}
\label{tab:nb_variants}
\resizebox{\textwidth}{!}{%
\begin{tabular}{@{}lcccc@{}}
\toprule
\textbf{Model} & \textbf{Accuracy} & \textbf{Precision} & \textbf{Recall} & \textbf{F1 Score} \\
\midrule
Gaussian NB     & 0.8100 & 0.8102 & 0.8076 & 0.8084 \\
Multinomial NB  & 0.7033 & 0.7726 & 0.7192 & 0.6922 \\
Bernoulli NB    & 0.4767 & 0.6356 & 0.5089 & 0.3438 \\
\bottomrule
\end{tabular}%
}
\end{table}

\subsubsection*{KNN Performance for Different k Values}
\begin{table}[H]
\centering
\caption{KNN Performance for Different k Values}
\label{tab:knn_k_values}
\begin{tabular}{@{}lcccc@{}}
\toprule
\textbf{k} & \textbf{Accuracy} & \textbf{Precision} & \textbf{Recall} & \textbf{F1 Score} \\
\midrule
1 & 0.7683 & 0.7616 & 0.7592 & 0.7603 \\
3 & 0.7591 & 0.7524 & 0.7482 & 0.7499 \\
5 & 0.7622 & 0.7572 & 0.7475 & 0.7508 \\
7 & 0.7561 & 0.7528 & 0.7381 & 0.7423 \\
9 & 0.7348 & 0.7322 & 0.7123 & 0.7166 \\
\bottomrule
\end{tabular}
\end{table}

\subsubsection*{KNN Comparison: KDTree vs BallTree}
\begin{table}[H]
\centering
\caption{KNN Comparison: KDTree vs BallTree}
\label{tab:knn_tree_comparison}
\begin{tabular}{@{}lcc@{}}
\toprule
\textbf{Metric} & \textbf{KDTree} & \textbf{BallTree} \\
\midrule
Accuracy & 0.7622 & 0.7622 \\
Precision & 0.7572 & 0.7572 \\
Recall & 0.7475 & 0.7475 \\
F1 Score & 0.7508 & 0.7508 \\
Training Time (s) & 0.0149 & 0.0092 \\
\bottomrule
\end{tabular}
\end{table}

\subsubsection*{SVM Performance with Different Kernels and Parameters}
\begin{table}[H]
\centering
\caption{SVM Performance with Different Kernels and Parameters}
\label{tab:svm_kernels}
\resizebox{\textwidth}{!}{%
\begin{tabular}{@{}lcccc@{}}
\toprule
\textbf{Kernel} & \textbf{Hyperparameters} & \textbf{Accuracy} & \textbf{F1 Score} & \textbf{Training Time (s)} \\
\midrule
Linear & C = 1.0 & 0.9268 & 0.9231 & 354.61 \\
Polynomial & C = 1.0, degree = 3, gamma = 'scale' & 0.6524 & 0.5248 & 1.20 \\
RBF & C = 1.0, gamma = 'scale' & 0.7165 & 0.6607 & 0.65 \\
Sigmoid & C = 1.0, gamma = 'scale' & 0.5366 & 0.5170 & 0.65 \\
\bottomrule
\end{tabular}%
}
\end{table}
\begin{table}[H]
\centering
\caption{Average 5-Fold Cross-Validation Results}
\label{tab:cv_comparison}
\small
\begin{tabular}{@{}lcccc@{}}
\toprule
\textbf{Model} & \textbf{Accuracy} & \textbf{Precision} & \textbf{Recall} & \textbf{F1-Score} \\
\midrule
Logistic Regression & 0.918 & 0.925 & 0.908 & 0.915 \\
GaussianNB & 0.831 & 0.758 & 0.981 & 0.855 \\
MultinomialNB & 0.768 & 0.755 & 0.771 & 0.762 \\
BernoulliNB & 0.879 & 0.872 & 0.888 & 0.880 \\
K-Neighbors Classifier & 0.791 & 0.829 & 0.719 & 0.770 \\
SVC (Linear) & 0.921 & 0.926 & 0.911 & 0.918 \\
SVC (Poly) & 0.651 & 0.960 & 0.379 & 0.544 \\
SVC (RBF) & 0.692 & 0.968 & 0.451 & 0.616 \\

\bottomrule
\end{tabular}
\end{table}

\begin{table}[H]
\centering
\caption{Final Test Set Results}
\label{tab:test_comparison}
\small
\begin{tabular}{@{}lcccc@{}}
\toprule
\textbf{Model} & \textbf{Accuracy} & \textbf{Precision} & \textbf{Recall} & \textbf{F1-Score} \\
\midrule
Logistic Regression & 0.9207 & 0.9287 & 0.9098 & 0.9167 \\
GaussianNB & 0.8353 & 0.7607 & 0.9855 & 0.8589 \\
MultinomialNB & 0.7713 & 0.7584 & 0.7753 & 0.7667 \\
BernoulliNB & 0.8810 & 0.8750 & 0.8913 & 0.8830 \\
K-Neighbors Classifier & 0.7957 & 0.8333 & 0.7246 & 0.7751 \\
SVC (Linear) & 0.9237 & 0.9292 & 0.9130 & 0.9210 \\
SVC (Poly) & 0.6554 & 0.9636 & 0.3840 & 0.5492 \\
SVC (RBF) & 0.6951 & 0.9708 & 0.4565 & 0.6212 \\

\bottomrule
\end{tabular}
\end{table}

\begin{table}[H]
\centering
\caption{K-Fold Cross-Validation Scores for Each Model (K=5)}
\label{tab:kfold_cv}
\begin{tabular}{@{}lccc@{}}
\toprule
\textbf{Fold} & \textbf{Na¨ıve Bayes Accuracy} & \textbf{KNN Accuracy} & \textbf{SVM Accuracy} \\
\midrule
Fold 1 & 0.8627 & 0.7643 & 0.9268 \\
Fold 2 & 0.8993 & 0.7712 & 0.9268 \\
Fold 3 & 0.8673 & 0.7666 & 0.9108 \\
Fold 4 & 0.8810 & 0.7529 & 0.9176 \\
Fold 5 & 0.8947 & 0.8009 & 0.9268 \\
\midrule
\textbf{Average} & \textbf{0.879} & \textbf{0.791} & \textbf{0.921} \\
\bottomrule
\multicolumn{4}{l}{}
\end{tabular}
\end{table}

\section{Best Practices Followed}
To ensure the experiment was methodologically sound and the results were reliable, several key machine learning best practices were followed:

\begin{itemize}
    \item \textbf{Robust Evaluation:} 5-Fold Cross-Validation was used instead of a single train-test split to get a more stable and accurate measure of each model's performance.

    \item \textbf{Comprehensive Metrics:} Models were judged on a suite of metrics (Accuracy, Precision, Recall, and F1-Score) to provide a holistic view of their performance.

    \item \textbf{Reproducibility:} The entire experiment was made reproducible by setting a consistent \texttt{random\_state} for data splits and model training.

    \item \textbf{Preventing Data Leakage:} Preprocessing steps like feature scaling were fitted only on the training data to prevent the model from gaining unfair knowledge of the test set, ensuring the results are realistic.
    
    \item \textbf{Baseline Modeling:} Simple models like Logistic Regression were used to establish a performance benchmark, which helps to justify the added value of more complex models.
\end{itemize}

\section{Conclusion}
\begin{itemize}
    \item \textbf{Champion Model Identified:} The Support Vector Classifier (SVC) with a Linear Kernel emerged as the top-performing model. It achieved the highest F1-Score of 0.9210 and an accuracy of 0.9237 on the final test set.

    \item \textbf{Linear Models Excelled:} The strong performance of the linear SVC, along with Logistic Regression, suggests the features in this dataset are largely linearly separable. More complex, non-linear models failed to provide any additional predictive advantage.

    \item \textbf{Ineffective Models:} K-Nearest Neighbors and SVMs with non-linear kernels (Polynomial, RBF, and Sigmoid) proved to be less effective for this specific problem. These models yielded significantly lower performance scores across all metrics.

    \item \textbf{Robust and Generalizable Results:} The chosen SVC (Linear) model demonstrated strong generalization, as its cross-validation F1-score (0.918) was very close to its final test set score. This indicates the model is robust and not overfitted to the training data.
\end{itemize}


\section{Learning Outcomes}
\begin{itemize}
    \item Gained experience implementing a full machine learning pipeline for a classification problem.
    \item Understood the importance of establishing a performance baseline with simple models.
    \item Learned to justify model selection based on dataset characteristics.
    \item Demonstrated the application of various classification algorithms and evaluation metrics.
\end{itemize}

\section*{GitHub Repository}
The complete source code for this experiment is available on GitHub:
\href{https://github.com/yuvashreeph/Machine-Learning/tree/main/Assignment3}{GitHub}

\end{document}
```